\chapter{Physical origins}

String amplitudes, gauge theory partition functions, and conformal
blocks each produce bar-cobar structures independently: the
generating function of $n$-point functions is the bar complex, and
amplitude factorization under degeneration is the coproduct.
This chapter traces the bar-cobar adjunction to its physical sources
and formulates each physics interface as a comparison problem with
named input data.

\section{4d/2d correspondence algebras}

\subsection{The 4d/2d correspondence (BLLPRR)}

\begin{construction}[4d/2d correspondence (BLLPRR)]
Starting from a 4d $\mathcal{N} = 2$ superconformal field theory $\mathcal{T}_{4d}$:
\begin{enumerate}[label=(\roman*)]
\item Select the \emph{Schur supercharge} $\mathbb{Q} = Q_1^- + \widetilde{S}^{1\dot{-}}$, a specific nilpotent supercharge in the $\mathcal{N}=2$ superconformal algebra ($\mathbb{Q}^2 = 0$).
\item Pass to the cohomology $H^*(\mathbb{Q})$ of local operators inserted at a fixed plane $\mathbb{R}^2 \subset \mathbb{R}^4$.
\item The resulting algebra $\mathcal{A}[\mathcal{T}_{4d}] = H^*(\mathbb{Q}, \text{local ops})$ carries the structure of a 2d vertex algebra (chiral algebra).
\end{enumerate}

This construction, due to Beem--Lemos--Liendo--Peelaers--Rastelli--van~Rees~\cite{beem-et-al}, produces vertex algebras from 4d SCFTs, by a mechanism distinct from the $\Omega$-background/Nekrasov partition-function
approach of the AGT correspondence below.
\end{construction}

\begin{theorem}[4d/2d is \texorpdfstring{$\Einf$}{E-infinity}-chiral]
The 2d chiral algebra $\mathcal{A}[\mathcal{T}_{4d}]$ obtained from the 4d/2d
correspondence is an $\Einf$-chiral algebra (vertex algebra).

The OPE of the resulting chiral algebra is generically non-commutative: the 4d/2d image of $SU(2)$ $\mathcal{N}=4$ SYM is the small $\mathcal{N}=4$ superconformal vertex algebra at $c = -9$, containing $\widehat{\mathfrak{sl}}_2$ at admissible level $k = -3/2$ (determined by $c = 3k/(k+2)$) as its R-symmetry current subalgebra. The $\Einf$ structure arises from the full vertex algebra axioms (locality, associativity of OPE), not from commutativity.
\end{theorem}

\begin{example}[Class \texorpdfstring{$\mathcal{S}$}{S}]
For class $\mathcal{S}$ theories of type $A_{n-1}$ on a Riemann surface $C$:
\[
\mathcal{A}[\mathcal{T}_{A_{n-1}}(C)] = \mathcal{W}^{-n}(\mathfrak{sl}_n)
\]
the W-algebra associated to $\mathfrak{sl}_n$ at the critical level $k = -n = -h^\vee$. At this special level the Sugawara construction degenerates and $\mathcal{W}^{-h^\vee}(\mathfrak{sl}_n)$ coincides with the Feigin--Frenkel center $\mathfrak{z}(\widehat{\mathfrak{g}})$, a commutative vertex algebra (see Feigin--Frenkel \cite{FF} for the center isomorphism and Arakawa \cite{Ara15} for related rationality results). The dependence on the Riemann surface $C$ and its pants decomposition appears at the level of conformal blocks and modules, not the algebra itself.
\end{example}


\section{Non-commutative Chern--Simons theory}
\subsection{Chern--Simons as source of chiral algebras}

\begin{construction}[CS boundary algebra]
3d Chern--Simons theory with gauge group $G$ at level $k$ on $M^3 = \Sigma \times \bR_+$
produces:
\begin{enumerate}[label=(\roman*)]
\item On the boundary $\Sigma$: WZW model at level $k$.
\item The chiral algebra is the affine Kac--Moody algebra $\hat{\mathfrak{g}}_k$.
\end{enumerate}
\end{construction}

\begin{remark}[Global de Rham obstruction; licensing tag $\epsilon$ via $\hypHTGlobalDR$]
\label{rem:ht-global-de-rham-physical-origins}
The physical-origin constructions of this chapter (CS-on-$\Sigma\times\R_+$
boundary WZW, Beem--Lemos 4d/2d Schur cohomology, AGT, twistor-string
holography, open-string boundary VOAs) all read formal-local
holomorphic-topological data and license a chiral algebra on a
compact bulk source. Each such promotion is conditional on
$\hypHTGlobalDR$: vanishing $H^1(X,\C_X) = 0$ for the holomorphic
de~Rham complex of the global manifold, descent of the local
Hamiltonian sheaf~\cite[\S4--5]{CG17}, satisfaction of the quantum
master equation, vanishing of the anomaly polynomial, and
Costello--Gwilliam locality. The local model is a section of the
obstruction sheaf; the global passage requires the obstruction
class to vanish
(\texttt{mixed-holomorphic-topological-strings:main.tex:3200--3210,
3232--3233}). For 3d CS on $\Sigma\times\R_+$ the obstruction
vanishes automatically because $\R_+$ is contractible and $\Sigma$
is one-complex-dimensional; for 4d/2d, AGT, twistor and 6d hCS the
obstruction is non-trivial and the Costello selection rule
(Proposition~\ref{prop:ht-twistor-deligne}) is its maximal solution
on the quartic locus.
\end{remark}

\begin{remark}[Non-commutative Chern--Simons comparison problem]
Noncommutative Chern--Simons theory on $\bR^3_\theta$ produces
$\Eone$-chiral algebras:
\begin{enumerate}[label=(\roman*)]
\item The boundary theory is a non-commutative WZW model.
\item The OPE is R-twisted with $R$ depending on the non-commutativity parameter $\theta$.
\item At $\theta = 0$, we recover the standard $\Einf$-chiral Kac--Moody.
\end{enumerate}
\end{remark}

\begin{remark}[Seiberg--Witten input]
The $R$-twisted OPE structure is consistent with Definition~\ref*{V1-def:chiral-ass-operad}; deriving it from non-commutative CS requires the Seiberg--Witten deformation quantization map~\cite{SW99}.
\end{remark}


\subsection{\texorpdfstring{Noncommutative geometry and associative / topological $\mathsf{E}_1$ structures}{Noncommutative geometry and associative / topological E1 structures}}

\begin{construction}[NC torus as associative deformation]
\index{noncommutative torus!topological $E_1$ structure@topological $\mathsf{E}_1$ structure}
The noncommutative torus $T^2_\theta$ ($\theta \in \mathbb{R} \setminus \mathbb{Q}$), with Moyal star product $f \star_\theta g$ deforming $\{x,y\}=\theta$, carries a topological $\Eone$-algebra structure ($VU = e^{2\pi i\theta}UV$, associative but noncommutative). This is \emph{not} a Beilinson--Drinfeld $\Eone$-chiral algebra on a curve.
\end{construction}

\begin{remark}[Koszul duality as Morita equivalence]
\index{Morita equivalence!Koszul duality interpretation}
The noncommutative-torus comparison asks for a topological
$\Eone$-Koszul duality model whose duality map exchanges the deformation
parameter:
\[
\mathcal{A}_\theta^! \;\cong\; \mathcal{A}_{1/\theta}.
\]
This coincides with Rieffel's Morita equivalence $T^2_\theta \sim_{\mathrm{Morita}} T^2_{1/\theta}$
for irrational $\theta$, reinterpreted as $\Eone$-Koszul duality.
The bar complex $\barB(\mathcal{A}_\theta)$ computes the cyclic homology
$\mathrm{HC}_*(\mathcal{A}_\theta)$, which by Connes' theorem equals the de~Rham
cohomology $H^*_{\mathrm{dR}}(T^2)$, independent of $\theta$ (topological invariance
of cyclic homology for smooth deformations).
\end{remark}

\begin{remark}[Morita input]
Rieffel Morita equivalence and Connes' $\mathrm{HC}_*(A_\theta)$ computation are established. The comparison consists of constructing the topological $\Eone$ bar complex for $\mathcal{A}_\theta$ and proving that its cohomology recovers the Koszul dual at $1/\theta$.
\end{remark}


\section{Gauge theory and D-branes}
\subsection{D-brane vertex algebras}

\begin{construction}[Open string VOA]
Open strings on D-branes in type IIB produce boundary chiral algebras determined by the brane configuration.
\end{construction}

\begin{remark}[D-brane algebras are \texorpdfstring{$\Eone$}{E1}]
Open string vertex algebras on non-trivial D-brane configurations are generically $\Eone$-chiral: Chan--Paton factors give matrix-valued OPE ($\Phi^{ij}(z)\Psi^{k\ell}(w) \sim \delta^{jk}\cdots$), breaking $\Einf$-commutativity.
\end{remark}

\begin{remark}[String-field-theory input]
Matrix-valued OPE from Chan--Paton factors is compatible with $\Eone$-chiral structure (Definition~\ref*{V1-def:chiral-ass-operad}); deriving it from open string field theory requires the cubic string-field-theory vertex of~\cite{Zwi93}.
\end{remark}


\subsection{\texorpdfstring{D-brane categories from $\Eone$-module Koszul duality}{D-brane categories from E1-module Koszul duality}}

\begin{remark}[HMS from \texorpdfstring{$\Eone$}{E1}-module Koszul duality]
\index{homological mirror symmetry!Koszul duality interpretation}
\index{D-brane!module Koszul duality}
In string theory, D-branes on a target space $X$ form the derived category
$D^b(\mathrm{Coh}(X))$ (B-branes) or the Fukaya category $\mathrm{Fuk}(X)$
(A-branes). For a chiral algebra $\mathcal{A}$ on a curve $\Sigma$, the
module category $\mathrm{Mod}(\mathcal{A})$ is the mathematical incarnation
of the category of branes.

Under $\Eone$-Koszul duality, the module bar-cobar functor
$\Phi\colon \mathrm{Mod}(\mathcal{A}) \xrightarrow{\sim}
\mathrm{CoMod}(\mathcal{A}^!)$
of Theorem~\ref*{V1-thm:e1-module-koszul-duality} exchanges module and comodule
categories. When $X$ is a Calabi--Yau surface and $\mathcal{A}$ is the chiral
algebra of the sigma model on $X$, the comparison identifies this
functor with homological mirror symmetry (Kontsevich):
\[
D^b\bigl(\mathrm{Coh}(X)\bigr)
\;\simeq\;
D^b\bigl(\mathrm{Fuk}(X^\vee)\bigr)
\]
for a mirror pair $(X, X^\vee)$. In this interpretation, the bar-cobar
adjunction is the proposed passage from B-branes to A-branes: the bar construction
$\barB(\mathcal{A})$ computes the Ext-algebra of the diagonal brane,
and the cobar construction recovers the endomorphism algebra of the dual
collection on $X^\vee$. The required input is the identification of
sigma-model chiral algebra modules with the A- and B-brane categories.
\end{remark}

\begin{remark}[Brane-category input]
Theorem~\ref*{V1-thm:e1-module-koszul-duality} stands independently; the HMS interpretation requires identifying chiral algebra modules with brane categories (boundary CFT input) and passing from the abstract functor $\Phi$ to derived category equivalences (tilting generation).
\end{remark}


\section{AGT correspondence connections}

\begin{theorem}[AGT correspondence]
(Alday--Gaiotto--Tachikawa) For 4d $\mathcal{N}=2$ $SU(2)$ gauge theory with $N_f=4$: $Z_{\mathrm{Nekrasov}}(\epsilon_1,\epsilon_2,a;q) = \langle V_{\alpha_1}(z_1)\cdots V_{\alpha_4}(z_4)\rangle_{\mathrm{Liouville}}$, with $c = 1+6Q^2$, $Q=b+b^{-1}$, $b^2=\epsilon_1/\epsilon_2$.
\end{theorem}

\begin{corollary}[Virasoro from gauge theory]
The Virasoro algebra ($\Einf$-chiral) arises from the AGT computation; $c$ is determined by $\Omega$-background parameters.
\end{corollary}


\begin{remark}[\texorpdfstring{$q$}{q}-AGT comparison problem]
The 5d lift of AGT (adding a circle) gives:
\[
Z_{\mathrm{5d Nekrasov}}(q, t) = \text{(conformal blocks of } \mathcal{W}_{q,t})
\]
where $\mathcal{W}_{q,t}$ is the quantum W-algebra, an $\Eone$-chiral algebra.
\end{remark}

\begin{remark}[Five-dimensional gauge-theory input]
$\mathcal{W}_{q,t}$ as an $\Eone$-chiral algebra is well-defined
(Frenkel--Reshetikhin). The equality with 5d Nekrasov partition
functions~\cite{AY10} is the remaining 5d gauge-theoretic comparison.
\end{remark}
